\begin{figure}[htbp]
\centering
\fbox{\parbox{0.9\columnwidth}{\centering
\textbf{Global Vulnerability Exploitation Statistics (2023-2025)}\\
\vspace{0.5cm}
\textbf{768 CVEs} exploited in 2024 (up from 639 in 2023) \cite{vulncheck2024trends}\\
\textbf{23.6\%} zero-day exploits \cite{infosec2024kevs}\\
\textbf{14\%} of breaches via vulnerability exploitation \cite{verizon2024dbir}\\
\textbf{180\%} increase vs 2023 \cite{globenews2024dbir}\\
\textbf{159 KEVs} in Q1 2025 \cite{vulncheck2024trends}\\
\vspace{0.3cm}
\rule{0.8\linewidth}{0.5pt}\\
\textbf{PRM Protection:} 11 PROG Handlers covering 13 Attack Vectors
}}
\caption{Global vulnerability exploitation statistics (2023-2025) showing the threat landscape that PRM's PROG handlers are designed to address. Data reflects publicly reported exploitation events across all sectors \cite{vulncheck2024trends,verizon2024dbir}.}
\label{fig:vuln_stats}
\end{figure}

\textbf{Note}: These data reflect publicly reported vulnerability-exploitation events globally (all sectors). Because public breach reporting aggregates many exploit types under 'vulnerability exploitation' and generally does not reveal whether the target was an HPC/data-center, the statistics here should be considered a lower bound on exposure. Finer-grained breakdowns (container escape, kernel rootkits, memory-code injection, etc.) are not systematically reported; their prevalence in HPC environments remains unknown.